% Style based on the supplied Waseem_Generic_AI_ML_GenAI_AgenticAI.tex.
% Python replaces template tokens. Edit layout here, content in JSON/portfolio JS.
\documentclass[9pt,letterpaper]{extarticle}
\usepackage[margin=0.55in]{geometry}
\usepackage[T1]{fontenc}
\input{glyphtounicode}
\pdfgentounicode=1
\usepackage[utf8]{inputenc}
\usepackage{lmodern}
\usepackage{microtype}
\usepackage{pifont}
\usepackage{xcolor}
\usepackage{enumitem}
\usepackage{titlesec}
\usepackage[hidelinks]{hyperref}
\usepackage{tabularx,array,parskip,ragged2e,needspace}
\pagestyle{empty}
\setlength{\parindent}{0pt}
\setlength{\parskip}{1pt}
\definecolor{accent}{HTML}{155B8A}
\titleformat{\section}{\large\bfseries\color{accent}}{}{0pt}{}[\vspace{0.5pt}\color{accent}\titlerule]
\titlespacing*{\section}{0pt}{3pt}{2pt}
\setlist[itemize]{leftmargin=1.15em,itemsep=0.8pt,topsep=1pt,parsep=0pt,partopsep=0pt}
\newcommand{\role}[4]{%
  \Needspace{4\baselineskip}\textbf{#1}\hfill\textbf{#2}\\[-1pt]
  \textit{#3}\hfill\textit{#4}\par}
\newcommand{\entry}[2]{\Needspace{4\baselineskip}\textbf{#1}\hfill\textit{#2}\par}
\newcommand{\skillrow}[2]{\textbf{#1:} & #2\\}
\hypersetup{pdftitle={Waseem Raza, PhD - Faculty \& Academic},pdfauthor={Waseem Raza, PhD}}
\begin{document}
\begin{center}
{\LARGE\bfseries Waseem Raza, PhD}\\[2pt]
{\large\bfseries Researcher \& Engineering Educator \textbar{} AI, Wireless Systems \& Applied Machine Learning}\\[1pt]
{\small AI for Wireless Networks \textbar{} Reliable Learning Systems \textbar{} Wearable Intelligence \textbar{} Engineering Education}\\[3pt]
{\footnotesize New Jersey, United States \quad \ding{37}\enspace \href{tel:+17328039703}{+1 732-803-9703} \quad \ding{41}\enspace \href{mailto:waseem.placements@gmail.com}{waseem.placements@gmail.com} \textbar{} \href{mailto:waseemraza844@gmail.com}{waseemraza844@gmail.com}\\[1pt]\href{https://www.linkedin.com/in/waseem-raza-phd-9a737858/}{LinkedIn} \textbar{} \href{https://github.com/WaseemRaza844}{GitHub} \textbar{} \href{https://scholar.google.com/citations?user=d58Drr0AAAAJ\&hl=en}{Google Scholar} \textbar{} \href{https://waseemraza844.github.io/portfolio/}{Portfolio}}
\end{center}
\vspace{-5pt}
\Needspace{3\baselineskip}
\section{Professional Summary}
Researcher and engineering educator with a PhD in Electrical and Computer Engineering and 5+ years of research and teaching experience spanning AI-enabled wireless networks, resilient machine learning, propagation modeling, optimization, 5G/6G systems, and wearable intelligence. Research combines analytical modeling, Python-based simulation, experimental validation, deep learning, transfer learning, domain adaptation, generative modeling, and reinforcement learning. Teaching and mentoring experience includes communications, wireless systems, programming, laboratory instruction, student projects, research methods, technical writing, and interdisciplinary collaboration.
\Needspace{3\baselineskip}
\section{Technical Skills}
\Needspace{2\baselineskip}\noindent{\small\textbf{Research Areas:} AI for wireless networks, resilient and autonomous networks, propagation modeling, digital twins, RAN optimization, wearable intelligence, and data-efficient learning}\par
\Needspace{2\baselineskip}\noindent{\small\textbf{Machine Learning:} Deep learning, transfer learning, domain adaptation, adversarial learning, generative modeling, reinforcement learning, classification, time-series analysis, and model evaluation}\par
\Needspace{2\baselineskip}\noindent{\small\textbf{Wireless Systems:} LTE/4G, 5G-NR, NTN, mobility, radio access networks, propagation, outage management, network measurements, and service assurance}\par
\Needspace{2\baselineskip}\noindent{\small\textbf{Research Methods:} Experimental design, simulation, statistical analysis, cross-validation, robustness evaluation, reproducible workflows, and technical communication}\par
\Needspace{2\baselineskip}\noindent{\small\textbf{Engineering \& Data:} Python, SQL, PyTorch, TensorFlow, scikit-learn, XGBoost, Pandas, NumPy, cloud analytics, data pipelines, and visualization}\par
\Needspace{2\baselineskip}\noindent{\small\textbf{Teaching \& Mentoring:} Wireless communications, digital communications, networking, signals and systems, machine learning, programming, experimental methods, and student research development}\par
\Needspace{4\baselineskip}\textbf{Leadership and Professional Strengths}\par
\begin{itemize}
\item \textbf{Research leadership:} Develop research questions into reproducible modeling, simulation, experimentation, validation, and publication workflows.
\item \textbf{Teaching and mentoring:} Connect engineering fundamentals with applied examples and support students in technical problem solving, implementation, experimentation, and communication.
\item \textbf{Industry translation:} Bridge academic methods with practical telecommunications problems through experience in university research and industry R\&D environments.
\item \textbf{Interdisciplinary collaboration:} Work across wireless communications, machine learning, data science, and wearable sensing while communicating results to technical and domain stakeholders.
\item \textbf{Persistence and rigor:} Approach difficult research problems through systematic experimentation, careful validation, iterative refinement, and attention to reproducibility.
\item \textbf{Continuous development:} Extend teaching and research capability through ongoing study in generative AI, agentic systems, cloud engineering, and emerging wireless technologies.
\end{itemize}
\Needspace{3\baselineskip}
\section{Professional Experience}
\role{Machine Learning Engineer}{02/2025 - Present}{AT\&T, contractor via Innovatix Inc.}{Bedminster, NJ}
\begin{itemize}\interlinepenalty=10000
\item Translate machine-learning research into operational telecommunications systems, emphasizing domain adaptation, reliability, model validation, and responsible decision support.
\item Develop reusable cross-domain evaluation and error-analysis methods for AI systems operating under distribution shift and heterogeneous operational conditions.
\item Collaborate across engineering, software, and public-safety teams to connect research methods with real-world network-assurance requirements.
\end{itemize}
\role{Advanced-Degrees Intern}{06/2024 - 08/2024}{AT\&T Labs}{Austin, TX}
\begin{itemize}\interlinepenalty=10000
\item Analyzed 500K+ LTE/5G-NR diagnostic logs and applied ML methods to detect anomalies, reducing diagnostic time by 20\%.
\item Investigated 5G NR-NTN cell-selection behavior and improved decision algorithms by 15\% through data-driven analysis and validation.
\item Conducted structured analysis of large operational datasets and communicated technical findings to engineering teams supporting next-generation wireless systems.
\end{itemize}
\role{Advanced Wireless Technology Development Intern}{09/2022 - 01/2023}{Futurewei Technologies}{Schaumburg, IL}
\begin{itemize}\interlinepenalty=10000
\item Improved autoencoder-based channel-state-information compression by 25\% using transfer learning and systematic experimentation on large wireless datasets.
\item Conducted 50+ experiments to validate model performance, robustness, and generalization under changing data conditions.
\item Collaborated with research and engineering stakeholders to translate modeling results into recommendations for next-generation communication systems.
\end{itemize}
\role{Advanced-Degrees Intern}{06/2022 - 08/2022}{AT\&T Labs}{Bedminster, NJ}
\begin{itemize}\interlinepenalty=10000
\item Designed ML and statistical methods for analyzing users' timing in 5G standalone networks, improving analytical precision by 20\%.
\item Analyzed large-scale user trace data with SQL and visualization tools and communicated findings to senior technical leadership.
\item Converted raw operational traces into interpretable performance indicators to support diagnosis and network optimization.
\end{itemize}
\role{Graduate Research Assistant}{09/2019 - 07/2025}{AI4Networks, University of Oklahoma}{Tulsa, OK}
\begin{itemize}\interlinepenalty=10000
\item Designed reusable Python simulation and data-generation platforms for large-scale experimentation, including a 3GPP-compliant network simulator and an AI-enabled 5G testbed.
\item Built transfer-learning, adversarial-learning, GAN, autoencoder, and reinforcement-learning pipelines for sparse-data and cross-domain modeling.
\item Developed robustness evaluation frameworks for models operating under distribution shifts, missing data, and limited labeled samples; led 100+ experiments from problem formulation through validation.
\item Developed scalable self-healing and outage-management frameworks that improved modeled network resilience by 30\%.
\item Published peer-reviewed research, presented technical results internationally, and mentored junior researchers in ML implementation, experimental design, and technical writing.
\end{itemize}
\role{Lecturer}{}{Engineering and Technology Programs}{}
\begin{itemize}\interlinepenalty=10000
\item Taught and mentored undergraduate engineering students in communications, programming, and foundational engineering subjects.
\end{itemize}
\Needspace{3\baselineskip}
\section{Education}
\entry{PhD, Electrical and Computer Engineering}{09/2019 - 07/2025}
University of Oklahoma, Norman, Oklahoma, USA\par
\entry{MSc, Telecommunication Engineering}{08/2014 - 08/2016}
University of Engineering and Technology, Taxila, Pakistan\par
\entry{BSc, Telecommunication Engineering}{11/2010 - 06/2014}
University of Engineering and Technology, Taxila, Pakistan\par
\Needspace{3\baselineskip}
\section{Selected Publications}
\entry{Real-World Medication Adherence Monitoring Using Hybrid Neural Networks and Smartwatches}{2026}
\textit{HealthINF \textbar{} Accepted}\par
\entry{On Multi-Parameter Optimization and Proactive Reliability in Cellular Networks}{2025}
\textit{Journal article}\par
\entry{An AI-Driven Framework for Enhancing Resilience in Propagation Models to Enable Digital Twin}{2024}
\textit{IEEE PIMRC}\par
\entry{Data-Driven Intelligent Outage Management for High Shadowing Environments in 5G\&B Networks}{2024}
\textit{IEEE SmartNets}\par
\entry{\href{https://arxiv.org/abs/2203.15899}{Toward a Hybrid RF/Optical Lunar Communication System (LunarComm)}}{2022}
\textit{IEEE Network}\par
\entry{Towards Positioning Error Impact Characterization and Minimization in User-Centric RAN}{2022}
\textit{IEEE WCNC}\par
\Needspace{3\baselineskip}
\section{Selected Projects}
\entry{SyntheticNET: 3GPP-Compliant Network Simulation}{2020-2024}
\textit{Research Project}\par
\begin{itemize}\interlinepenalty=10000
\item Python-based simulation environment for realistic mobility, propagation, and data-driven experimentation in cellular networks.
\end{itemize}
{\small\textit{Tools \& methods:} Python, 3GPP, Simulation, Wireless Networks\par}
\entry{TurboRAN Experimental Platform}{2020-2024}
\textit{Research Infrastructure}\par
\begin{itemize}\interlinepenalty=10000
\item Multi-band, multi-tier experimental platform supporting evaluation of emerging AI-enabled radio-access-network solutions.
\end{itemize}
{\small\textit{Tools \& methods:} 5G Testbeds, RAN, Experiment Design, Network Measurement\par}
\entry{Wearable Intelligence for Health Monitoring}{2024-2026}
\textit{Research Project}\par
\begin{itemize}\interlinepenalty=10000
\item Deep-learning methods for continuous recognition of medication and dietary activities from wrist-worn inertial sensors.
\end{itemize}
{\small\textit{Tools \& methods:} Wearable AI, Time Series, Deep Learning, Smartwatch Sensing\par}
\Needspace{3\baselineskip}
\section{Certifications}
\Needspace{3\baselineskip}\noindent\textbf{1. Machine Learning} {\small --- Stanford University / Coursera \textbar{} Completed}\par
\textbf{Focus:} Core supervised and unsupervised learning, model evaluation, optimization, and practical ML system development.\par
{\small\textbf{Tools \& Skills:} Supervised Learning, Unsupervised Learning, Model Evaluation, Optimization\par}
\vspace{2pt}
\Needspace{3\baselineskip}\noindent\textbf{2. AI for Telecommunications Specialization} {\small --- \textcolor{blue}{\underline{\href{https://coursera.org/verify/specialization/D8Y6MVC0XM51}{AI CERTs / Coursera}}} \textbar{} Completed (3/3) courses \textbar{} 06/2026 \textbar{} \textcolor{blue}{\underline{\href{https://waseemraza844.github.io/portfolio/certificates/wireless/ai-for-telecom.pdf}{Certificate}}}}\par
\textbf{Focus:} Three-course specialization connecting AI foundations with telecommunications network optimization, security, and customer experience.\par
{\small\textbf{Tools \& Skills:} AI for Telecommunications, Network Optimization, Cybersecurity, Customer Experience, Internet of Things, Python, TensorFlow\par}
\vspace{2pt}
\Needspace{3\baselineskip}\noindent\textbf{3. 4G Network Fundamentals} {\small --- \textcolor{blue}{\underline{\href{https://coursera.org/verify/DLW1YN474VZJ}{Institut Mines-Télécom / Coursera}}} \textbar{} Completed (1/1) courses \textbar{} 08/2026 \textbar{} \textcolor{blue}{\underline{\href{https://waseemraza844.github.io/portfolio/certificates/wireless/4g-network-fundamentals.pdf}{Certificate}}}}\par
\textbf{Focus:} Foundational study of fourth-generation cellular networks.\par
{\small\textbf{Tools \& Skills:} 4G Networks, Cellular Networks, Telecommunications\par}
\vspace{2pt}
\Needspace{3\baselineskip}\noindent\textbf{4. Generative AI Fundamentals Specialization} {\small --- \textcolor{blue}{\underline{\href{https://coursera.org/verify/specialization/IPQKDOHF3IV2}{IBM / Coursera}}} \textbar{} Completed (5/5) courses \textbar{} 07/2026 \textbar{} \textcolor{blue}{\underline{\href{https://waseemraza844.github.io/portfolio/certificates/genai/ibm-genai-fundamentals/ibm-genai-fundamentals.pdf}{Certificate}}}}\par
\textbf{Focus:} Five-course foundation in generative AI, prompt engineering, LLMs, and practical enterprise applications.\par
{\small\textbf{Tools \& Skills:} Generative AI, Prompt Engineering, Foundation Models, Responsible AI, IBM watsonx, Hugging Face, Business Applications\par}
\vspace{2pt}
\Needspace{3\baselineskip}\noindent\textbf{5. AWS Generative AI and AI Agents with Amazon Bedrock} {\small --- \textcolor{blue}{\underline{\href{https://coursera.org/verify/professional-cert/XMNOUZNOW9M0}{AWS/Coursera}}} \textbar{} Completed (3/3) courses \textbar{} 08/2026 \textbar{} \textcolor{blue}{\underline{\href{https://waseemraza844.github.io/portfolio/certificates/genai/aws-bedrock/aws-generative-ai-agents-amazon-bedrock.pdf}{Certificate}}}}\par
\textbf{Focus:} Three-course professional certificate focused on developing, customizing, optimizing, and automating generative AI applications and intelligent agents with Amazon Bedrock.\par
{\small\textbf{Tools \& Skills:} Amazon Bedrock, Generative AI, AI Agents, Foundation Models, Knowledge Bases, Retrieval-Augmented Generation, LangChain, Amazon Q Developer, Model Customization, Prompt Engineering, Natural Language Processing, Text Generation, Summarization, Generative AI Application Development\par}
\vspace{2pt}
\end{document}
